\documentclass[11pt]{article}
\usepackage[letterpaper,margin=1in]{geometry}
\usepackage[T1]{fontenc}
\usepackage{lmodern}
\usepackage{amsmath,amssymb,amsthm}
\usepackage{xurl}
\usepackage[hidelinks]{hyperref}
\usepackage{microtype}
\setlength{\parindent}{0pt}
\setlength{\parskip}{0.55em}
\newcommand{\Z}{\mathbb Z}
\newcommand{\N}{\mathbb N}
\newtheorem*{theorem}{Theorem}
\title{Disproof of Conjecture 00000001735\\[0.35em]
\large An attained maximum cannot have finitely many maximizing right-hand sides}
\author{C0ldSmi1e}
\date{October 9, 2026}
\begin{document}
\maketitle
\vspace{-1.2em}

\section{The statement and the counting convention}

Both versions of the original conjecture assert that the maximum solution count
for $X^3-2Y^3=k$ is $12$ and that only finitely many $k$ attain $12$.
They also promise an explicit finite list through an unspecified unit-group
torsion condition. We disprove the conjunction of the maximum and finiteness
assertions, which is a necessary part of the full statement.

We use the ordinary integer-solution interpretation of a Thue equation:
\[
 S_k=\{(x,y)\in\Z^2:x^3-2y^3=k\},\qquad N(k)=\#S_k,
 \qquad k\in\Z.
\]
Solutions are distinct ordered pairs. Signs and zero coordinates are included.
Neither language restricts solutions to primitive pairs, restricts $k$ to
cube-free integers, or identifies parameters up to multiplication by cubes.
None of those additional restrictions is imposed here. The customary alternative
of requiring $k\ne0$ is also treated below.

\begin{theorem}
If $N(k)\le12$ for every integer $k$ and $N(k_0)=12$ for some integer $k_0$,
then infinitely many integers $k$ have $N(k)=12$. Consequently the conjecture,
as stated, is false.
\end{theorem}

\begin{proof}
First $S_0=\{(0,0)\}$. Indeed, $x^3=2y^3$ with $y\ne0$ would also have
$x\ne0$. Taking absolute values and comparing exponents of the prime $2$
in the two positive integers gives
\[
 3v_2(|x|)=1+3v_2(|y|),
\]
which is impossible modulo $3$. Thus $y=0$ and then $x=0$.
In particular $N(0)=1$, so a parameter $k_0$ with $N(k_0)=12$ is nonzero.

For every nonzero integer $t$, homogeneity gives an injective map
\[
 S_{k_0}\longrightarrow S_{t^3k_0},\qquad (x,y)\longmapsto(tx,ty),
\]
because $(tx)^3-2(ty)^3=t^3(x^3-2y^3)$ and multiplication by $t$ is
injective on integers. Hence $N(t^3k_0)\ge12$. The assumed global upper bound
forces $N(t^3k_0)=12$.

Take $t=n+1$ for $n\in\N$. The parameters $(n+1)^3k_0$ are pairwise distinct:
one may cancel the nonzero $k_0$, and cubing integers is injective.
They form an infinite subset of the claimed finite set of maximizers.
\end{proof}

\newpage
\section{Scope of the disproof}

The conclusion is the incompatibility of the two assertions. We do not determine
the actual numerical maximum, prove that the number $12$ is attained, or compute
an actual list of maximizing parameters.

For the parameter domain $\Z\setminus\{0\}$, an upper bound of $12$ extends
to all integers because $N(0)=1$. Excluding zero does not change the set of
twelve-solution parameters. This variant is separately formalized and refuted.

The undefined torsion-condition phrase is not given an invented meaning.
Whatever additional enumeration assertion it expresses, the full source still
implies the maximum-and-finiteness core refuted above. In particular a finite
explicit list cannot repair that contradiction.

\section{Formalization and reproduction}

The Lean project uses Lean 4.19.0 and Mathlib v4.19.0, with dependency revisions
pinned in its manifest. In \nolinkurl{lean/Solution.lean},
\nolinkurl{Conjecture1735.solutionSet} is exactly the displayed subset of
$\Z\times\Z$ and \nolinkurl{solutionCount} is
\nolinkurl{Cardinal.mk (solutionSet k)}. Using cardinal numbers counts the whole
set even before its finiteness is known; it never assigns an artificial finite
count to an infinite set. A count equal to the finite cardinal $12$ is the
ordinary statement that there are exactly twelve solutions.

\nolinkurl{MaximumIsTwelve} means a universal upper bound of twelve together
with actual attainment. \nolinkurl{Conjecture} conjoins it with
\nolinkurl{Set.Finite} of the full set of attaining parameters. The main endpoints
are:
\begin{itemize}
 \item \nolinkurl{Conjecture1735.maximum_twelve_implies_infinite_maximizers}:
 the asserted maximum implies \nolinkurl{Set.Infinite} of the maximizer set;
 \item \nolinkurl{Conjecture1735.conjecture_false}:
 the unconditional negation of \nolinkurl{Conjecture};
 \item \nolinkurl{Conjecture1735.nonzero_conjecture_false}:
 the unconditional negation of the nonzero-parameter variant.
\end{itemize}

The zero-fiber argument, scaling map, injection, cardinal inequality, and infinite
range are all proved in Lean. No search bound, numerical experiment, auxiliary
mathematical program, or unproved finiteness theorem is used. The submitted
inspection file expands the definitions and prints endpoint types and axiom
dependencies. The portable verifier rebuilds the project, replays the inspection
and full compiled-declaration audit with warnings treated as errors, and checks
the source, dependency, and declaration records. Its accompanying controls test
the rejection of invalid or altered verification inputs. Exact commands and
results are recorded in \nolinkurl{README.md} and \nolinkurl{VERIFICATION.md}.

The original bilingual source is included as \nolinkurl{ORIGINAL.md}, with
its exact hash recorded in the verification manifest.
The author derived the scaling argument from that source before an accidental
exposure to unrelated task-status summaries; the access record discloses that
event. A fresh reviewer independently identified the same argument before
reading the proof. Full provenance and local verification records accompany
the submission; local validation is distinct from maintainer acceptance.

\end{document}
